\section{Where GeFL-F loses the rare classes}\label{ch:diagnosis}

The paper finds a candidate loss on each of the first three routes of \cref{sec:s-routes}: the optimiser in Stage II,
the last layer in Stage III, and the synthetic data in Stage III. Each subsection below gives the intuition with the
running example, the precise statement, and what the runs show. The strength of the evidence differs between the
three, and the paper says so.

\subsection{The generator: coupled weight decay under Adam}\label{sec:s-decay}

\paragraph{The situation} In one round of Stage II, every client starts from the global generator, takes some local
Adam steps on its own features, and returns its copy; the server averages the copies. Consider the class row $r_9$ of
the running example on client 2, which holds no image of class 9. Client 2 never generates a class-9 feature during
training, so its data loss does not depend on $r_9$ at all. The only force on $r_9$ is the weight decay. With coupled
decay the penalty $\frac\lambda2\|\theta\|^2$ is part of the loss, so the gradient of client 2's loss with respect to
$r_9$ is exactly
\[
g=\lambda r_9 .
\]
With plain SGD this would move $r_9$ by $\eta\lambda r_9$ per step, a tiny amount: with the generator's learning rate
$\eta=10^{-3}$ and decay $\lambda=10^{-3}$, one part in a million of the row per step. Adam behaves very differently.

\paragraph{Why Adam makes the decay large} Adam divides each gradient coordinate by a running estimate of its own
magnitude. With the running means $m$ of the gradient and $v$ of its square (decay rates $\beta_1=0.9$,
$\beta_2=0.999$) and their bias-corrected versions $\hat m,\hat v$, a step is
\[
\theta\leftarrow\theta-\eta\,\frac{\hat m}{\sqrt{\hat v}+\epsilon}.
\]
Take one coordinate of $r_9$ equal to $0.5$, so $g=\lambda\cdot0.5=5\times10^{-4}$. On the first step $m=0.1\,g$ and
$v=0.001\,g^2$; the bias corrections divide them by $1-0.9$ and $1-0.999$, so $\hat m=g$ and $\hat v=g^2$. The step is
$\eta\,g/(|g|+\epsilon)\approx\eta=10^{-3}$, in the direction of the sign of the coordinate. Now repeat with a decay
100 times smaller, $\lambda=10^{-5}$: then $g=5\times10^{-6}$, still much larger than $\epsilon=10^{-8}$, and the step is
again $\approx10^{-3}$. Adam's step does not depend on the size of the gradient, only on its sign and consistency. A
client that lacks class 9 therefore moves every coordinate of $r_9$ toward zero by about one learning rate per step,
whatever the decay, 2,000 times faster than SGD would in this example.

\begin{lemma*}[coupled decay under Adam]
On a client without class $c$, the gradient of the loss with coupled L2 regularisation $\lambda$ with respect to $r_c$
is exactly $\lambda r_c$. Since Adam is invariant to a positive rescaling of the gradient sequence up to its $\epsilon$,
each step of such a client moves every coordinate of $r_c$ toward zero by about the learning rate $\eta$, independently
of $\lambda$, as long as $\lambda|r_{c,j}|\gg\epsilon$. With $\lambda=0$ this shrinkage vanishes; with decoupled weight
decay (AdamW) it becomes a factor $1-\eta\lambda$ per step, which is negligible here.
\end{lemma*}

Two remarks complete the picture. A client keeps its Adam state between rounds while its parameters are reset to the
global ones, so a single step can be somewhat larger than $\eta$. And the effect is specific to \emph{federated}
training: on pooled data, every batch that contains class 9 gives $r_9$ a data gradient, Adam's estimate $v$ of the
squared gradient is dominated by these, and the decay stays small. This holds as long as class-9 batches come within
the memory of $v$, about $1/(1-\beta_2)=1000$ steps.

\paragraph{The holders do not compensate} One might expect the clients that hold class 9 to push $r_9$ back. They
push it toward a good value, not toward a larger norm: their data gradient fits the row and need not lengthen it. They
also see class 9 rarely. Stage II trains the generator for 5 local epochs per round in batches of 64 (from the code).
Client 1 of the running example holds 4 class-9 images among 913, so at most 4 of its 15 batches per epoch, and at most
20 of its 75 steps per round, contain the class. On every other step its gradient on $r_9$ is again only the decay
$\lambda r_9$, and Adam shrinks the row exactly as on a client without the class. So most holders shrink a rare class's
row most of the time.

\paragraph{Averaging turns this into contraction} After the round, the server averages the returned rows. Let
$\Delta_k=\|r_{k,9}\|-\|r_9\|$ be the change of the row's length on client $k$, $\bar G$ the mean change over the
holders, $\bar D$ the mean shrinkage over the non-holders, and $p=m/K$ the share of clients that hold the class.

\begin{lemma*}[contraction under flat averaging]
Under $r_c^+=\frac1K\sum_k r_{k,c}$, $\|r_c^+\|\le\|r_c\|+p_c\bar G_c-(1-p_c)\bar D_c$, with equality if and only if all
$r_{k,c}$ are non-negative multiples of one another. The row contracts whenever $p_c<\bar D_c/(\bar G_c+\bar D_c)$.
\end{lemma*}

The proof is the triangle inequality: the length of an average is at most the average of the lengths. The
inequality is often strict, because clients that fit class 9 to their own few images move the row in different
directions, and averaging vectors that point in different directions shortens them further. For a feel of the
numbers: if eight holders lengthen the row by $0.02$ each and two non-holders shorten it by $0.05$, the bound allows a
gain of $0.8\cdot0.02-0.2\cdot0.05=0.006$. If the holders' average change is zero, which is plausible when they rarely
see the class, the row shrinks by at least $0.2\cdot0.05=0.01$ every round, over 200 rounds.

The paper also gives a predictor of the direction, counting Adam steps:
\[
\kappa_c=\frac{\sum_{k\in H_c}q_k f_{kc}}{\sum_{k\notin H_c}q_k},
\]
where $H_c$ is the set of holders, $q_k$ the number of local steps of client $k$, and $f_{kc}$ the probability that a
batch of client $k$ contains class $c$. The numerator counts the steps on which a holder actually sees the class, the
denominator the steps on which a non-holder decays it. When $\kappa_c<1$ the row drifts toward zero. The paper is
explicit that this is a heuristic, not a bound: it ignores the momentum Adam carries between rounds, and it was
\emph{not} checked against the measured rows.

\paragraph{The remedies} Two changes remove the effect. Decoupled weight decay (AdamW) applies the decay outside
Adam's normalisation, so the shrinkage becomes the negligible factor $1-\eta\lambda$. Or the rows can be averaged over
their holders only, weighted by how much of the class each holds (class-count weighted averaging, CCWA, in
\cref{sec:s-gen}): the non-holders' term $\bar D$ then drops out of the lemma.

\paragraph{What was measured} The paper's figure of the class rows (\cref{fig:s-rows}) shows that after Stage II the
rows of the rarest classes are several times shorter than those of the most frequent ones. That is the predicted
effect. It is not yet evidence that the shrinkage \emph{causes} the rare classes' low accuracy, and the paper says
so, for a reason that comes from a second measurement.

After Stage I the server trains a classifier, the \emph{referee}, on the features of its validation images, and asks
it to classify 300 generated features of each class. The share it assigns to the intended class is the generator's
\emph{fidelity} for that class. If the decay were what makes rare classes fail, the rare classes' fidelity would be
low and the frequent classes' high. Instead GeFL-F's generated CIFAR-10-LT features are recognised at
\R{main/cifar10/GeFL-F/fidelity_all} on average (chance is 10\%), \R{main/cifar10/GeFL-F/fidelity_head} on the head
classes and \R{main/cifar10/GeFL-F/fidelity_tail} on the tail classes. On CIFAR-100-LT the mean is
\R{main/cifar100/GeFL-F/fidelity_all}, chance level (\cref{tab:s-fidelity}). GeFL-F's generator carries little class
information for \emph{any} class. The row shrinkage is a measured property of the generator; whether it causes the
rare-class loss is open. The direct test is GeFL-F with AdamW and nothing else changed, one of the runs not yet
returned (\cref{ch:results}). The referee's accuracy on held-out real features was not logged, so the fidelity values
compare generators with each other, not with real data.

\begin{figure}[t]
\centering
\includegraphics[width=0.62\linewidth]{\genpath/figures/fig_rows.pdf}
\caption{The paper's figure of the class rows: length of each class row of GeFL-F's generator after Stage II, classes
from the most frequent (left) to the least frequent (right); CIFAR-10-LT, mean over seeds 3--5. Look at how the
length falls toward the rare classes.}
\label{fig:s-rows}
\end{figure}

\begin{table}[t]
\centering\small
\caption{The paper's fidelity table: the referee's per-class accuracy on generated features (\%), averaged over all
classes and over the head (30\% most frequent), medium and tail (40\% least frequent) classes; mean $\pm$ standard
deviation over seeds 3--5. FedDF uses GeFL-F's generator. FedECS-LT's rows belong to \cref{ch:method}.}
\label{tab:s-fidelity}
\slottab{fidelity}{}
\end{table}

\subsection{The last layer: logit corrections cannot add information}\label{sec:s-lastlayer}

\paragraph{The situation} In Stage III, client 2's head is trained on client 2's features only, in the real phase.
Its last layer therefore learns client 2's class proportions: classes it holds often get large biases, the classes
it lacks (class 9 among them) get pushed down by every cross-entropy step, because cross-entropy lowers the logits of
all wrong classes. The natural fix from long-tail learning is to correct the loss with the client's class
frequencies. \emph{Logit adjustment} trains with $\ell(f+\tau\log\hat\pi_k,\,y)$, where $\hat\pi_k$ is client $k$'s
class distribution (with add-one smoothing, so that absent classes get a small non-zero frequency) and $\tau=1$.
Balanced Softmax and FedLC are of the same form: a constant vector $a_k$ added to the logits during training.

\begin{lemma*}[logit corrections are reparametrisations]
Let a head have a free last-layer bias $b$, trained by SGD without decay on $b$, and let $a_k\in\Rset^C$ be a constant
of client $k$, as in logit adjustment ($a_k=\tau\log\pi_k$), FedLC or Balanced Softmax. Training with the loss
$\ell(f+a_k,y)$ from $(\theta_0,b_0)$ produces exactly the trajectory of training with $\ell(f,y)$ from
$(\theta_0,b_0+a_k)$: every parameter except $b$ is the same at every step. FedRS, which scales the logits of absent
classes by $\alpha\in(0,1]$, is cross-entropy in the rescaled last-layer variables $\alpha w_c,\alpha b_c$ of those
classes, trained with the smaller learning rate $\alpha^2\eta$.
\end{lemma*}

The proof: $f+a_k=We+(b+a_k)$, and moving from $b$ to $b+a_k$ is a translation, which does not change gradients. So
training with the corrected loss is plain training that starts from the bias $b_0+a_k$; at test time the correction is
not applied, which amounts to subtracting $a_k$ from the logits. This is exactly \emph{post-hoc logit adjustment} of a
head trained from a shifted starting bias.

\paragraph{A three-class example} A client holds 70 images of class A, 30 of class B and none of class C. With add-one
smoothing its frequencies are $(71,31,1)/103$, so $a=\log\hat\pi\approx(-0.37,\,-1.20,\,-4.63)$. At test time the
corrected head adds $+4.63$ to class C's logit relative to plain training, $+3.43$ more than to class B. That raises
class C for \emph{every} input by the same amount. It cannot make class C high for images of class C and low for others,
because the weight vector $w_C$ never received a single gradient that points toward class-C features: the client has
no class-C features. The correction moves a threshold; it adds no information about the class. The body of the head
can still learn a somewhat different representation, because the shifted logits change the gradients that reach it,
but nothing tells it what class C looks like. A head trained on one client is an extreme case of what the literature
calls minority collapse: for the classes the client lacks, the imbalance ratio is infinite.

\paragraph{What follows} A last layer that knows the absent classes needs information from other clients. The
partition-invariance lemma of \cref{ch:sums} says where to get it: a last layer that is a function of sums over all
clients equals the one computed on pooled data. The paper's answer is the \emph{exact readout} of \cref{sec:s-readout},
a classifier computed from summed embedding statistics.

\paragraph{What was measured} Two measurements support this part of the diagnosis, and they are the strongest
evidence in the paper's analysis. First, replacing GeFL-F's trained last layers by the exact readout, with nothing else
changed, raises its balanced accuracy by \R{main/cifar10/GeFL-F/gain_er} points on CIFAR-10-LT and
\R{main/cifar100/GeFL-F/gain_er} on CIFAR-100-LT. Most of what GeFL-F loses is in the last layer. Second, the lemma
predicts that Balanced Softmax, FedLC and FedRS differ only in their last layers. Once these are replaced by the same
readout, the three methods reach \R{main/cifar10/Balanced Softmax/er_bal}, \R{main/cifar10/FedLC/er_bal} and
\R{main/cifar10/FedRS/er_bal} on CIFAR-10-LT, and \R{main/cifar100/Balanced Softmax/er_bal},
\R{main/cifar100/FedLC/er_bal} and \R{main/cifar100/FedRS/er_bal} on CIFAR-100-LT: the same, as predicted.

\subsection{The synthetic data: label-free features penalise the logit spread}\label{sec:s-spread}

\paragraph{The situation} In the synthetic phase, labels are drawn uniformly over the $C$ classes and the generator
makes a feature for each. If the generated feature carries its label, the head learns the class. If it does not, as
the fidelity of \cref{tab:s-fidelity} suggests for GeFL-F, the label is effectively independent of the feature, and
the training target is the uniform distribution.

\begin{lemma*}[label-free synthetic data]
If the synthetic label $y$ is uniform on $\{1,\dots,C\}$ and independent of the generated feature, then for every logit
vector $f$,
\[
\E_y\,\ell(f,y)=\LSE(f)-\bar f=\log C+\KL(u\,\|\,\mathrm{softmax}f)=\log C+\frac1{2C}\sum_c(f_c-\bar f)^2+O(\|f-\bar f\|^3),
\]
where $\LSE(f)=\log\sum_c e^{f_c}$, $\bar f$ is the mean logit and $u$ the uniform distribution.
\end{lemma*}

In words: averaged over the random label, the cross-entropy is smallest when all logits are equal, and it grows
roughly with the variance of the logits. Training on label-free features therefore pulls the logits together. The
paper measures this as the \emph{logit spread}: the mean over images of the standard deviation of a head's logits
across classes.

A numerical check with $C=10$: equal logits give the loss $\log10=2.303$. The logits $(2,0,\dots,0)$ give
$\LSE=\log(e^2+9)=2.797$ and $\bar f=0.2$, so the loss is $2.597$. The quadratic term alone predicts
$2.303+\frac1{20}(1.8^2+9\cdot0.2^2)=2.483$; the rest is the higher-order term. The more confident the head, the larger
the penalty, whatever the class.

\paragraph{What survives the alternation} The synthetic phase is followed by the real phase, which pushes the logits
apart again for the classes the client holds. The paper analyses the alternation in a simplified model with quadratic
losses (an appendix lemma). The round converges to a point where, along every direction of the parameters that the
real loss leaves free, the synthetic loss alone decides. For a head trained on one client, the directions that
concern its absent classes are almost free in this sense, so on them the label-free synthetic phase decides: it
flattens exactly the logits the client most needs. Not exactly free, because cross-entropy on the real data still
pushes the absent classes' logits down on data the head separates. Logit adjustment shifts those targets.

\paragraph{What was measured} The paper's figure of the logit spread (\cref{fig:s-spread}) shows, in every Stage-III
round on CIFAR-10-LT, that the synthetic phase reduces the spread of the heads' logits on the clients' real features,
and that it reduces it less for FedECS-LT than for GeFL-F. Two caveats keep this from being a clean test. FedECS-LT
also differs from GeFL-F in its number of synthetic epochs, its soft targets and logit adjustment, so the comparison
does not isolate the labels. And FedDF, which trains on GeFL-F's own label-free generator, loses \emph{less} spread
than FedECS-LT. So the size of the spread change is not determined by whether the generated features carry their
label. The lemma is exact; how much it matters in the real pipeline is not settled.

\begin{figure}[t]
\centering
\includegraphics[width=0.62\linewidth]{\genpath/figures/fig_spread.pdf}
\caption{The paper's logit-spread figure: change of the logit spread during the synthetic phase of each Stage-III
round (after minus before; negative means the logits were pulled together), CIFAR-10-LT, mean over clients and seeds
3--5. Look at the levels: GeFL-F lowest, FedECS-LT in between, FedDF highest.}
\label{fig:s-spread}
\end{figure}

What would a method look like that uses exact statistics exactly where these losses occur?
